% V2原书第12章思考题12.1--12.10，扫描PDF40--41；与章末习题同号但独立记录。
\section*{第12章\quad 静电场：思考题}
\paragraph{思考题 12.1}
点电荷公式在距离 (r\to0) 时电场趋于无穷大，怎样理解点电荷模型的适用范围？

\paragraph{思考题 12.2}
从两点电荷间静电力与电荷乘积成正比说明库仑定律与电场叠加原理的关系及其局限。

\paragraph{思考题 12.3}
说明电场的操作性定义（试验电荷所受力除以试验电荷量）与孤立点电荷的库仑电场公式之间的联系；试验电荷应满足什么条件？

\paragraph{思考题 12.4}
电场线与电通量、场强大小和电荷的正负有何关系？

\paragraph{思考题 12.5}
把三个相同正电荷置于正三角形三顶点，为什么不能把三角形中心取作球心、直接用高斯定律算该处电场大小？实际中心电场是多少？

\paragraph{思考题 12.6}
某闭合曲面S的电通量为零，是否能断定面上所有位置的电场也为零？

\paragraph{思考题 12.7}
若闭合曲面每一点电场均为零，可否断定曲面内部一定没有任何净电荷或点电荷？

\paragraph{思考题 12.8}
电场线能否在无自由电荷的空间中突然中断？根据电场通量与散度作答。

\paragraph{思考题 12.9}
高斯定律与库仑定律之间在静电学中有什么联系？为何高斯定律并非任意场对称性下求E的捷径？

\paragraph{思考题 12.10}
真空两块大而近似无限平行板各带+q与-q，面积S、间距d：三种不同电场/力计算为什么可能给出矛盾值？推导正确的理想平板间相互作用力。
